\chapter{AC Steady State: Impedance, Admittance, and Power}
\label{ch:ac}
\Cref{ch:foundations} introduced the components---\(R\), \(L\), \(C\), and
sources---and \cref{ch:complex} introduced the phasor as the complex number that
stands in for a steady sinusoid. This chapter puts the two together. In the
\textbf{AC steady state}---a circuit driven by a sinusoid long enough for
transients to die out---every voltage and every current is a sinusoid at the
source frequency, so each can be written as a phasor. The element laws, rewritten
in phasors, then give the central circuit quantities: impedance, admittance, and
complex power. This is the phasor machinery of \cref{ch:complex} applied to actual
circuits.

\section{From Element Laws to Impedance}
Recall the main ideas of \cref{ch:complex}: for a phasor quantity, \(\dd/\dd t\)
becomes multiplication by \(\jj\omega\). Apply that to each element law in turn,
and watch the calculus disappear.

\emph{Resistor.} Ohm's law \(v=Ri\) has no derivative, so it passes through
unchanged: \(V=R\,I\).

\emph{Capacitor.} The element law is \(i=C\,\dd v/\dd t\). Substituting the phasor
forms \(v(t)=\sqrt2\,\Re\{V\ee^{\jj\omega t}\}\) and
\(i(t)=\sqrt2\,\Re\{I\ee^{\jj\omega t}\}\), the derivative brings out a factor
\(\jj\omega\):
\[
I\,\ee^{\jj\omega t}=C\frac{\dd}{\dd t}\bigl(V\ee^{\jj\omega t}\bigr)
=\jj\omega C\,V\,\ee^{\jj\omega t}
\quad\Longrightarrow\quad
I=\jj\omega C\,V .
\]
\emph{Inductor.} Likewise \(v=L\,\dd i/\dd t\) becomes \(V=\jj\omega L\,I\).

In every case the phasor voltage is a \emph{complex multiple} of the phasor
current---exactly the algebraic form of Ohm's law, with a complex proportionality
constant. Solving each relation for the ratio \(V/I\) defines that constant. It is
the \textbf{impedance}:
\[
Z=\frac{V}{I}=R+\jj X,
\]
a frequency-dependent generalization of resistance. For the capacitor, dividing
\(I=\jj\omega C V\) through gives
\[
Z_C=\frac{V}{I}=\frac{1}{\jj\omega C},
\]
which is where the \(1/\jj\omega\) comes from: the capacitor law puts \(\jj\omega\)
on the \emph{current} side, so it lands in the denominator of \(V/I\). Multiplying
top and bottom by \(-\jj\) (since \(1/\jj=-\jj\)) puts it in rectangular form,
\(Z_C=-\jj/(\omega C)\). Its real part \(R\) is
resistance (dissipation); its imaginary part \(X\) is \textbf{reactance} (energy
storage, no dissipation). For the three elements,
\[
Z_R=R,\qquad
Z_L=\jj\omega L,\qquad
Z_C=\frac{1}{\jj\omega C}=-\jj\frac{1}{\omega C}.
\]
An inductor contributes positive reactance \(X_L=\omega L\); a capacitor
contributes negative reactance, of magnitude \(X_C=1/(\omega C)\), so that
\(Z_C=-\jj X_C\). At low frequency a capacitor's impedance is large
(it blocks DC) and an inductor's is small (it passes DC); at high frequency the
roles reverse.

\begin{physicalbox}
\textbf{A chemical-engineering analogy.} Circuits obey the same
across/through (effort/flow) structure as the hydraulic and thermal systems of
process dynamics:
\begin{center}
\begin{tabularx}{0.92\textwidth}{@{}L{0.26\textwidth}Y@{}}
\toprule
Electrical & Process / hydraulic analog \\
\midrule
Voltage \(V\) (across) & Pressure or head (driving force) \\
Current \(I\) (through) & Volumetric flow rate \\
Resistance \(R\) & Flow resistance (pressure drop per unit flow) \\
Capacitance \(C\) & Tank capacitance (volume stored per unit head) \\
Inductance \(L\) & Fluid inertance (inertia opposing flow change) \\
Impedance \(Z\) & Frequency-domain ``flow resistance'' \\
\bottomrule
\end{tabularx}
\end{center}
So \emph{impedance} is the generalized flow resistance of a network in the
frequency domain, and \emph{reactance} is the piece contributed by storage---a
surge tank or fluid inertia opposes a \emph{change} in flow without dissipating
energy, storing it and handing it back a quarter cycle later. That energy return,
not loss, is exactly why reactance shifts phase by \(90^\circ\) instead of
consuming power.
\end{physicalbox}

\section{Series Impedance: Where the RLC Impedance Comes From}
Impedances in series add, exactly like resistors, because the same current flows
through each and the voltages sum (KVL). For \(R\), \(L\), \(C\) in series that
gives the impedance directly:
\[
Z=Z_R+Z_L+Z_C=R+\jj\omega L+\frac{1}{\jj\omega C}
=R+\jj\left(\omega L-\frac{1}{\omega C}\right).
\]
The reactance is the \emph{difference} \(\omega L-1/(\omega C)\); it vanishes at
resonance, a fact \cref{ch:rlc} builds on.

\begin{figure}[htbp]
\centering
\begin{circuitikz}
\draw (0,0) to[sV,l=\(v_s\)] (0,2)
      to[R,l=\(R\)] (2.6,2)
      to[L,l=\(L\)] (5.2,2)
      to[C,l=\(C\)] (5.2,0)
      -- (0,0);
\end{circuitikz}
\caption{A series RLC circuit whose impedance
\(Z=R+\jj(\omega L-1/\omega C)\) is just the sum of the three element
impedances.}
\end{figure}

\section{Parallel Impedance}
The dual arrangement puts the three elements across one node pair
(\cref{fig:rlc-parallel-z}), so they share a \emph{voltage} instead of a current.
Impedances no longer add---but their reciprocals do, because KCL sums the branch
currents:
\[
I=I_R+I_L+I_C
=\frac{V}{R}+\frac{V}{\jj\omega L}+\jj\omega C\,V
\quad\Longrightarrow\quad
\frac{I}{V}=\frac{1}{R}+\jj\left(\omega C-\frac{1}{\omega L}\right).
\]
Note the sign structure: for the parallel circuit the \emph{difference} is
\(\omega C-1/(\omega L)\), the mirror of the series circuit's
\(\omega L-1/(\omega C)\). It too vanishes at \(\omega_0=1/\sqrt{LC}\)---but there
it makes \(I/V\) a \emph{minimum} rather than \(V/I\), so a parallel resonator
presents a high impedance where a series resonator presents a low one
(\cref{ch:rlcpar}).

Rather than invert this expression every time, it is cleaner to give the reciprocal
of impedance its own name, which the next section does.

\begin{figure}[htbp]
\centering
\begin{circuitikz}[scale=1]
\draw (0,0) to[sV,l=\(v_s\)] (0,2.2);
\draw (0,2.2) -- (5,2.2);
\draw (0,0) -- (5,0);
\draw (1.8,2.2) to[R,l=\(R\)] (1.8,0);
\draw (3.4,2.2) to[L,l=\(L\)] (3.4,0);
\draw (5,2.2) to[C,l=\(C\)] (5,0);
\node[ground] at (2.6,0) {};
\end{circuitikz}
\caption{A parallel RLC circuit. All three elements share the same voltage, so the
branch currents add (KCL) and it is the \emph{reciprocals} of the impedances that
sum: \(I/V=1/R+\jj(\omega C-1/\omega L)\).}
\label{fig:rlc-parallel-z}
\end{figure}

\section{Admittance and Susceptance}
For parallel networks it is easier to work with the reciprocal of impedance, the
\textbf{admittance}
\[
Y=\frac{1}{Z}=G+\jj B,
\]
whose real part \(G\) is conductance and imaginary part \(B\) is
\textbf{susceptance}. Just as series impedances add, \emph{parallel admittances
add}, because the same voltage appears across each branch and the currents sum
(KCL):
\[
Y_{\mathrm{eq}}=\sum_k Y_k.
\]
In the process analogy, admittance is ``flow per unit driving force''---a
conductance to flow---and parallel pipes add their flow conductances the same way.

\begin{workedbox}[: a parallel branch]
A \(\SI{100}{\ohm}\) resistor is in parallel with a capacitor whose reactance is
\(\SI{100}{\ohm}\) at the frequency of interest (\cref{fig:rc-parallel-ex}). Their
admittances are
\[
Y_R=\frac{1}{100}=\SI{0.010}{\siemens},
\qquad
Y_C=\frac{1}{-\jj100}=\jj\,\SI{0.010}{\siemens},
\]
so \(Y_{\mathrm{eq}}=0.010+\jj0.010=\SI{0.0141}{\siemens}\angle45^\circ\), and the
parallel impedance is
\[
Z_{\mathrm{eq}}=\frac{1}{Y_{\mathrm{eq}}}\approx\SI{70.7}{\ohm}\angle{-45^\circ}
=50-\jj50\ \Omega.
\]
Adding admittances turned a parallel combination into one addition. Note that the
\(45^\circ\) admittance angle became \(-45^\circ\) on inversion, and that neither
\(50\) nor \(-50\) is obvious from the \(\SI{100}{\ohm}\) parts---parallel
combinations are where working in \(Y\) earns its keep.
\end{workedbox}

\begin{figure}[htbp]
\centering
\begin{circuitikz}[scale=1]
\draw (0,0) to[short,-o] (0,0);
\draw (0,2) node[left]{\(\circ\)} -- (2.8,2);
\draw (0,0) node[left]{\(\circ\)} -- (2.8,0);
\draw (1.2,2) to[R,l={\(R=\SI{100}{\ohm}\)}] (1.2,0);
\draw (2.8,2) to[C,l={\(X_C=\SI{100}{\ohm}\)}] (2.8,0);
\node at (-0.75,1) {\(Z_{\mathrm{eq}}\)};
\end{circuitikz}
\caption{The parallel \(R\)--\(C\) branch of the worked example, seen as a
one-port of impedance \(Z_{\mathrm{eq}}\).}
\label{fig:rc-parallel-ex}
\end{figure}

\section{Phasor Lead and Lag}
The \(90^\circ\) phase relationships are not extra facts to memorize---they fall
directly out of the element differential equations, and the derivative is doing all
the work.

Take a capacitor with \(v(t)=V_p\cos\omega t\). Its element law gives the current by
differentiating:
\[
i=C\frac{\dd v}{\dd t}=-\omega C V_p\sin\omega t
=\omega C V_p\cos\bigl(\omega t+90^\circ\bigr),
\]
using \(-\sin\theta=\cos(\theta+90^\circ)\). The current's cosine peaks a quarter
cycle \emph{earlier} than the voltage's: in a capacitor, \textbf{current leads
voltage by \(90^\circ\)}. The same conclusion in phasor form is immediate---
multiplying by \(\jj\omega\) adds \(90^\circ\), because \(\jj=1\angle90^\circ\):
\[
I=\jj\omega C\,V=\omega C\,V\angle{+90^\circ}.
\]
An inductor is the dual. With \(i(t)=I_p\cos\omega t\),
\(v=L\,\dd i/\dd t\) puts the \(+90^\circ\) on the \emph{voltage} instead, so
\textbf{voltage leads current by \(90^\circ\)}. In both cases the derivative in the
element law is the phase shift, and \(\jj\omega\) is that derivative wearing a
phasor's clothing.

Consequently a positive impedance angle means net inductive behavior and a negative
angle means net capacitive behavior---the angle simply reports which store wins.

\begin{workedbox}[: a lagging (inductive) load]
A \(\SI{50}{\ohm}\) resistor in series with \(X_L=\SI{50}{\ohm}\) gives
\(Z=50+\jj50=\SI{70.7}{\ohm}\angle{+45^\circ}\). For \(V=100\angle0^\circ\)~V,
\[
I=\frac{V}{Z}=\frac{100\angle0^\circ}{70.7\angle45^\circ}
=1.41\angle{-45^\circ}\ \si{\ampere}.
\]
The current angle is \emph{negative}: the current lags the voltage by
\(45^\circ\)---not the full \(90^\circ\), because the resistor contributes no phase
shift and pulls the total back toward zero. Positive \(Z\) angle, lagging current,
inductive.
\end{workedbox}

\begin{workedbox}[: a leading (capacitive) load]
Swap the inductor for a capacitor of the same reactance magnitude:
\(Z=50-\jj50=\SI{70.7}{\ohm}\angle{-45^\circ}\). Now
\[
I=\frac{100\angle0^\circ}{70.7\angle{-45^\circ}}
=1.41\angle{+45^\circ}\ \si{\ampere},
\]
so the current \emph{leads} by \(45^\circ\). Same magnitudes, opposite sign, and
the sign is the entire physical difference between the two loads.
\end{workedbox}

\begin{exambox}
If a question says ``lead'' or ``lag,'' sketch a tiny phasor diagram before
answering. Most mistakes come from recalling the right \(90^\circ\) relationship
with the direction reversed. Mnemonic: ELI the ICE man---in an inductor (L)
voltage E leads current I; in a capacitor (C) current I leads voltage E.
\end{exambox}

\section{Complex Power}
With RMS phasors and the passive sign convention, complex power uses the
conjugate of the current:
\[
S=V I^{*}=P+\jj Q .
\]
The real part \(P\) is average power actually dissipated (watts); the imaginary
part \(Q\) is reactive power (VAR, volt-amperes reactive), the energy sloshing
back and forth between the source and the reactive elements each cycle. The
conjugate is what makes \(P\) come out as the in-phase product, and the magnitude
\(|S|\) is the apparent power, measured in volt-amperes (VA).

\begin{workedbox}[: real and reactive power]
A load draws \(I=2\angle{-30^\circ}\ \si{\ampere}\) (RMS) from
\(V=120\angle0^\circ\ \si{\volt}\) (RMS)---current lagging, so the load is
inductive. Then
\[
S=VI^{*}=120\angle0^\circ\cdot 2\angle{+30^\circ}=240\angle30^\circ
=207.8+\jj120\ \mathrm{VA}.
\]
So \(P=\SI{207.8}{\watt}\) is dissipated and \(Q=120\,\mathrm{VAR}\) is exchanged with
the load's reactance. A purely reactive load (\(\pm90^\circ\)) would show
\(P=0\): it consumes no average power.
\end{workedbox}

\section{Worked Example: Series Reactance}
Find the impedance of a \(\SI{50}{\ohm}\) resistor in series with
\(\SI{10}{\micro\henry}\) at \(\SI{7.1}{\mega\hertz}\)
(\cref{fig:series-rl-ex})---a coil and its loss, seen at the low end of the
\SI{40}{\meter} band.

\begin{figure}[htbp]
\centering
\begin{circuitikz}[scale=1]
\draw (0,0) node[left]{\(\circ\)} to[R,l={\(R=\SI{50}{\ohm}\)}] (2.6,0)
      to[L,l={\(L=\SI{10}{\micro\henry}\)}] (5.4,0) node[right]{\(\circ\)};
\node at (2.7,-1.05) {\(Z=R+\jj\omega L\)};
\end{circuitikz}
\caption{The series \(R\)--\(L\) one-port of the worked example.}
\label{fig:series-rl-ex}
\end{figure}

First convert frequency to angular frequency, then take the reactance:
\[
\omega=2\pi(7.1\times10^6)=\SI{4.46e7}{\radian\per\second},
\qquad
X_L=\omega L=(4.46\times10^7)(10\times10^{-6})\approx\SI{446}{\ohm}.
\]
Thus \(Z=50+\jj446\ \Omega\), and in polar form
\[
|Z|=\sqrt{50^2+446^2}\approx\SI{449}{\ohm},
\qquad
\theta=\operatorname{atan2}(446,50)\approx83.6^\circ .
\]
The angle is close to \(+90^\circ\): the coil's reactance dwarfs its resistance, so
at this frequency the part behaves almost like a pure inductor, and the current will
lag the voltage by nearly a quarter cycle.
